\section{Conclusion}
\label{sec:conclusion}

We asked a narrow question: once a detector has done its imperfect best, what
does the removal stage actually deliver? The answer, across ten
recall--precision cells, four architecture--attack combinations, three seeds,
five removal families, and a second signal domain, is that it delivers what the
detector's recall allows and very little more. Deletion behaves as a step
function. At full recall it meets the success criterion; at $0.95$ recall the
attack is suppressed but the clean-accuracy cost has already broken the
criterion; at $0.90$ recall it leaves a model that is $\fset$-free by
construction and still $100\%$ attackable, at every seed and every setting we
tried. False positives are a second-order cost by comparison, visible only when
they are numerous enough to remove a large block of clean data.

The one method that degrades gracefully is repair followed by inoculation, and
the qualification matters more than the result. It tolerates $r{=}0.5$ because
its inoculation batch supplies the trigger, and a model trained to ignore that
trigger ignores it whether or not every poisoned sample was found. When we
replaced the known trigger with a synthesised one, the recovered patch had zero
pixel overlap with the true trigger and attack success stayed above $95\%$. The
recipe is therefore a ceiling on what trigger-aware repair can achieve against a
fixed attack, not a blind post-training defence. Calibrate-then-remove inherits
the same dependency and, in our sweep, could not find a threshold that separated
poison from clean data; amplify-and-detect did not raise the detectability of
missed poison above the promotion gate; and the apparent success of feature-space
ascent was an artefact of frozen batch-norm statistics.

Two further results bound the claims that removal papers can make. Durability
probes show suppression rather than erasure: the oracle-deleted model resists
membership inference and output-divergence tests yet returns to $97.7\%$ attack
success within ten relearning steps. And the recall cliff is not a CIFAR
artefact; on RadioML automatic modulation classification the minimum recall
$\gstar$ traces the same boundary as a function of SNR, rising where the clean
task itself becomes difficult.

The practical recommendation is a reporting one. A removal result quoted at a
single detector operating point cannot be interpreted, because the underlying
quantity is the amount of poison left behind, not the method's name or its
hyperparameters. Reporting removal as a curve over $(r,p)$, together with the
tolerance $\rstar$ and an explicit statement of whether the trigger was known,
would make results comparable and would expose the trigger dependence that our
experiments show is decisive. Under that reporting standard, the field's
post-training methods are bounded by detection, and improving the detector
remains the only lever that moves the bound.
